\section{Задачи поиска путей}

Одна из классических задач анализа графов~--- это задача поиска путей между вершинами с различными ограничениями.

При этом возможны различные постановки задачи.
С одной стороны, постановки различаются тем, что именно мы хотим получить в качестве результата.
Здесь наиболее частыми являются следующие варианты.

\begin{itemize}
    \item Наличие хотя бы одного пути, удовлетворяющего ограничениям, в графе.
          В данном случае не важно, между какими вершинами существует путь, важно лишь наличие его в графе.
    \item Наличие пути, удовлетворяющего ограничениям, между некоторыми вершинами: задача достижимости.
          При данной постановке задачи, нас интересует ответ на вопрос о достижимости вершины \circled{$v_i$} из вершины \circled{$v_j$} по пути, удовлетворяющему ограничениям.
          Такая постановка требует лишь проверить существование пути, но не его предоставления в явном виде.
    \item Поиск одного пути, удовлетворяющего ограничениям: необходимо не только установить факт наличия пути, но и  предъявить его.
          При этом часто подразумевается, что возвращается любой путь, являющийся решением, без каких-либо дополнительных ограничений.
          Хотя, например, в некоторых задачах дополнительное требование простоты или наименьшей длины выглядит достаточно естественным.
    \item Поиск всех путей: необходимо предоставить все пути, удовлетворяющие заданным ограничениям.
\end{itemize}

С другой стороны, задачи могут различаться ещё и тем, как фиксируются множества стартовых и конечных вершин.
Здесь возможны следующие варианты:
\begin{itemize}
    \item от одной вершины до всех,
    \item между всеми парами вершин,
    \item между фиксированной парой вершин,
    \item между двумя множествами вершин $V_1$ и $V_2$, что подразумевает решение задачи для всех $(v_i,v_j) \in V_1 \times V_2$.
\end{itemize}

Стоит отметить, что последний вариант является самым общим и остальные~--- лишь его частные случаи.
Однако этот вариант часто выделяют отдельно, подразумевая, что остальные варианты в него не включаются\sidenote{
    Часто это связано с тем, что для некоторых частных случаев можно построить специализированный алгоритм, который будет по каким-то показателям лучше, чем алгоритм для задачи общего вида.
}.

В итоге, перебирая возможные варианты желаемого результата и способы фиксации стартовых и финальных вершин, мы можем сформулировать достаточно большое количество задач.
Например, задачу поиска всех путей между двумя заданными вершинами, задачу поиска одного пути от фиксированной стартовой вершины до каждой вершины в графе, или задачу достижимости между всеми парами вершин.

Часто поиск путей сопровождается изучением их свойств, что далее приводит к формулированию дополнительных ограничений на пути в терминах этих свойств.
Например, можно потребовать, чтобы пути были простыми или не проходили через определённые вершины.
Один из естественных способов описывать свойства и, как следствие, задавать ограничения~--- это использовать ту алгебраическую структуру, из которой берутся веса рёбер графа%
\sidenote{На самом деле здесь наблюдается некоторая двойственность.
    С одной стороны, действительно, удобно считать, что свойства описываются в терминах некоторой заданной алгебраической структуры.
    Но, вместе с этим, структура подбирается исходя из решаемой задачи.}.

Предположим, что дан граф $\mbfscrG = \langle V, E, L\rangle $, где $L = (S, \oplus, \otimes)$~--- это полукольцо.
Тогда изучение свойств путей можно описать следующим образом:
\begin{equation}
    \label{eq:algPathProblem}
    \left\{\ (v_i, v_j, c) \mid c = \bigoplus_{v_i \pi_k v_j} \bigotimes_{(u, l, v) \in \pi } l \ \right\}.
\end{equation}

Иными словами, для каждой пары вершин, для которой существует хотя бы один путь, их соединяющий, мы агрегируем (с помощью операции $\oplus$ из полукольца) информацию обо всех путях между этими вершинами.
При этом информация о пути получается как свёртка меток рёбер пути с использованием операции $\otimes$\sidenote{Заметим, что детали свёртки вдоль пути зависят от свойств полукольца (и от решаемой задачи).
    Так, если полукольцо коммутативно, то нам не обязательно соблюдать порядок рёбер.
    В дальнейшем мы увидим, что данные особенности полукольца существенно влияют на особенности алгоритмов решения соответствующих задач.}.

Естественным требованием (хотя бы для прикладных задач, решаемых таким способом) является существование и конечность указанной суммы.
На данном этапе мы не будем касаться того, какие именно свойства полукольца могут нам обеспечить данное свойство, однако в дальнейшем будем считать, что оно выполняется.
Более того, будем стараться приводить частные для конкретной задачи рассуждения, показывающие, почему это свойство выполняется в рассматриваемых в задаче ограничениях.

Описанная выше задача общего вида называется анализом свойств путей алгебраическими методами (Algebraic Path Problem)~\sidecite{Baras2010PathPI} и предоставляет общий способ для решения широкого класса прикладных задач%
\sidenote{В работе \enquote{Path Problems in Networks}~\cite{Baras2010PathPI} собран действительно большой список прикладных задач с описанием соответствующих полуколец.
    Сводная таблица на страницах 58--59 содержит 29 различных прикладных задач и соответствующих полуколец.}.
Наиболее известными являются такие задачи, как построение транзитивного замыкания графа и поиск кратчайших путей (All Pairs Shortest Path или APSP).
Далее мы подробнее обсудим эти две задачи и предложим алгоритмы их решения.

